% Generated draft: requires mathematical and source review.
\documentclass[11pt]{article}
\usepackage[margin=25mm]{geometry}
\usepackage{fontspec}
\setmainfont{Latin Modern Roman}
\usepackage{amsmath,amssymb,mathtools,unicode-math}
\setmathfont{Latin Modern Math}
\usepackage{enumitem,needspace,xurl,hyperref}
\setlength{\parindent}{0pt}
\setlength{\parskip}{5pt}
\setlength{\emergencystretch}{3em}
\newcommand{\N}{\mathbb{N}}
\newcommand{\Z}{\mathbb{Z}}
\newcommand{\Q}{\mathbb{Q}}
\newcommand{\R}{\mathbb{R}}
\newcommand{\C}{\mathbb{C}}
\newcommand{\F}{\mathbb{F}}
\newcommand{\A}{\mathbb{A}}
\newcommand{\PP}{\mathbb{P}}
\newcommand{\E}{\mathbb{E}}
\newcommand{\eps}{\varepsilon}
\newcommand{\dd}{\,\mathrm{d}}
\begin{document}
\section*{UnsolvedMath ID: 30001903}
\textbf{Source problem code:} catalogue entry
\subsection{Definitions and mathematical statement}
The supplied record uses \(M\), the algebra \(CS(M)\), and a cyclic cocycle \(c\) called the Radul cocycle, but does not define them. In particular, it gives no hypotheses on \(M\), no specification of the algebra or its topology, and no formula or normalization for \(c\). These data remain unresolved; no particular interpretation of \(CS(M)\) is imposed here.

For the intended algebra \(A=CS(M)\), write \(HC^1(A)\) for its degree-one cyclic cohomology and \(HP^1(A)\) for its odd periodic cyclic cohomology, using the intended cohomology theory. Write
\[
\operatorname{per}\colon HC^1(A)\longrightarrow HP^1(A)
\]
for the canonical periodicization map in that theory. The algebraic or continuous cochain convention is not specified in the record.

\textbf{Requested identification.} With the intended hypotheses and definitions supplied, identify the class of the Radul cocycle and its periodic image:
\[
[c]\in HC^1(CS(M)),
\qquad
\operatorname{per}([c])\in HP^1(CS(M)).
\]
The record supplies neither a proposed identifying formula nor a target class, comparison map, or classification criterion beyond these cohomology groups. Consequently, this is a faithful statement of the requested identification, but the supplied information does not determine a fully specified mathematical problem.
\subsection{Short English statement}
Identify the cyclic-cohomology class represented by the Radul cocycle on the intended algebra \(CS(M)\), together with its image in odd periodic cyclic cohomology. The hypotheses on \(M\), the algebra, the cocycle, the cohomology convention, and the intended form of the identification require clarification. The supplied open status is a dated curation label.
\subsection{Sources}
\begin{itemize}
\item Local UnsolvedMath registry, permanent ID \(30001903\), source problem code \texttt{catalogue entry}, title ``Cyclic-Cohomology Classes of the Radul Cocycle,'' in the Oberwolfach Reports Open Problems collection. No individual proposer is supplied; the proposed year is recorded as 2011.
\item Supplied source citation: \emph{Noncommutative Geometry (2012)}, Oberwolfach Reports, DOI 10.4171/owr/2011/45, pp.\ 2549--2622. \url{https://doi.org/10.4171/owr/2011/45}. The background describes Oberwolfach Reports as documenting mathematical workshops and their open-problem discussions. The citation's 2012 label and the record's 2011 proposed year are retained without reconciliation.
\item Supplied literature source, labeled ``original OWR report'': \url{https://doi.org/10.4171/owr/2011/45}. The literature assessment was checked on \texttt{2026-08-21}, according to the supplied dated curation; this date and assessment have not been independently verified here. That assessment labels the problem open and says that the extracted question requests identification of a class associated with the Radul cyclic cocycle without specifying the target class or a complete classification result. The background expressly requests independent verification of time-sensitive status and cited claims.
\item The registry labels the statement \texttt{corrected\_verified} and reports that its mathematical content was checked against the official Oberwolfach report, with notation, context, and question form normalized without changing the source-level claim. This is a supplied verification assertion, not an independent source check performed for this draft.
\end{itemize}
\end{document}
